\textbf{Phần 1. Trắc nghiệm lựa chọn}
\setcounterpageref{ex}{0}

\begin{ex} %[3H1V18]
Một hình chữ nhật có 4 góc. Hỏi một hình chữ nhật có bao nhiêu góc vuông?
\choice
{$2$ góc vuông}
{\True $4$ góc vuông}
{$3$ góc vuông}
{$1$ góc vuông}
\loigiai{
Hình chữ nhật có 4 góc đỉnh đều là góc vuông.
\textbf{Đáp án của câu hỏi là B}
}
\end{ex}

\begin{ex} %[3H1V18]
Lúc đúng $3$ giờ, kim giờ và kim phút của đồng hồ tạo thành một góc gì?
\choice
{Góc nhọn}
{Góc tù}
{Góc bẹt}
{\True Góc vuông}
\loigiai{
Lúc 3 giờ đúng, kim phút chỉ số 12 và kim giờ chỉ số 3, hai kim tạo thành một góc vuông ($90^\circ$).
\textbf{Đáp án của câu hỏi là D}
}
\end{ex}

\begin{ex} %[3H1V18]
Trong các thời điểm sau, thời điểm nào mà hai kim đồng hồ (kim giờ và kim phút) tạo thành một góc vuông?
\choice
{\True $9$ giờ đúng}
{$6$ giờ đúng}
{$12$ giờ đúng}
{$4$ giờ đúng}
\loigiai{
Lúc 9 giờ đúng, kim phút chỉ số 12 và kim giờ chỉ số 9, hai kim vuông góc với nhau.
\textbf{Đáp án của câu hỏi là A}
}
\end{ex}

\begin{ex} %[3H1V18]
Một hình tứ giác có thể có nhiều nhất bao nhiêu góc vuông?
\choice
{$2$ góc vuông}
{$3$ góc vuông}
{\True $4$ góc vuông}
{$1$ góc vuông}
\loigiai{
Hình vuông và hình chữ nhật là các tứ giác có 4 góc vuông (nhiều nhất có thể).
\textbf{Đáp án của câu hỏi là C}
}
\end{ex}

\begin{ex} %[3H1V18]
Một hình tam giác có thể có nhiều nhất bao nhiêu góc vuông?
\choice
{$2$ góc vuông}
{\True $1$ góc vuông}
{$3$ góc vuông}
{Không có góc vuông nào}
\loigiai{
Một hình tam giác chỉ có thể có nhiều nhất 1 góc vuông (tam giác vuông).
\textbf{Đáp án của câu hỏi là B}
}
\end{ex}

\begin{ex} %[3H1V18]
Dụng cụ chuẩn thường dùng để kiểm tra một góc có phải là góc vuông hay không là
\choice
{Thước dây}
{Compa}
{Thước kẻ thẳng thông thường}
{\True Thước ê-ke}
\loigiai{
Thước ê-ke có góc vuông chuẩn dùng để kiểm tra góc vuông.
\textbf{Đáp án của câu hỏi là D}
}
\end{ex}

\begin{ex} %[3H1V18]
Lúc $2$ giờ đúng, góc tạo bởi kim giờ và kim phút của đồng hồ là góc gì?
\choice
{\True Góc nhọn (góc không vuông bé hơn góc vuông)}
{Góc vuông}
{Góc tù (góc không vuông lớn hơn góc vuông)}
{Góc bẹt}
\loigiai{
Lúc 2 giờ đúng, khoảng cách giữa hai kim là 2 khoảng giờ ($60^\circ$), bé hơn góc vuông nên là góc nhọn.
\textbf{Đáp án của câu hỏi là A}
}
\end{ex}

\begin{ex} %[3H1V18]
Lúc $4$ giờ đúng, góc tạo bởi kim giờ và kim phút của đồng hồ là
\choice
{Góc vuông}
{Góc không vuông bé hơn góc vuông (góc nhọn)}
{\True Góc không vuông lớn hơn góc vuông (góc tù)}
{Góc bẹt}
\loigiai{
Lúc 4 giờ đúng, góc giữa kim giờ và phút là 4 khoảng giờ ($120^\circ$), lớn hơn góc vuông nên là góc tù.
\textbf{Đáp án của câu hỏi là C}
}
\end{ex}

\begin{ex} %[3H1V18]
Bác thợ mộc ghép hai tấm gỗ hình chữ nhật lại với nhau để tạo thành một hình chữ nhật lớn hơn. Hỏi hình chữ nhật mới có tất cả bao nhiêu góc vuông?
\choice
{$8$ góc vuông}
{\True $4$ góc vuông}
{$6$ góc vuông}
{$2$ góc vuông}
\loigiai{
Bất kì hình chữ nhật nào cũng luôn có đúng 4 góc vuông ở 4 đỉnh.
\textbf{Đáp án của câu hỏi là B}
}
\end{ex}

\begin{ex} %[3H1V18]
Hình chữ nhật $ABCD$ có kẻ đường chéo $AC$. Hình vẽ khi đó xuất hiện bao nhiêu góc vuông?
\choice
{$2$ góc vuông}
{$6$ góc vuông}
{$8$ góc vuông}
{\True $4$ góc vuông}
\loigiai{
Hình chữ nhật có 4 góc vuông tại $A, B, C, D$. Đường chéo $AC$ chia góc $A$ và $C$ thành các góc nhọn, nhưng ở đỉnh $B$ và $D$ vẫn là góc vuông, còn góc toàn phần tại đỉnh $A$ và $C$ vẫn là góc vuông. Tổng cộng vẫn có 4 góc vuông.
\textbf{Đáp án của câu hỏi là D}
}
\end{ex}

\begin{ex} %[3H1V18]
Cho một góc vuông đỉnh $O$, cạnh $OA, OB$. Kẻ thêm một tia $OC$ nằm giữa hai cạnh $OA$ và $OB$. Khi đó trên hình có bao nhiêu góc nhọn?
\choice
{\True $2$ góc nhọn}
{$1$ góc nhọn}
{$3$ góc nhọn}
{$0$ góc nhọn}
\loigiai{
Góc ban đầu là góc vuông $\widehat{AOB}$. Tia $OC$ nằm trong chia thành hai góc $\widehat{AOC}$ và $\widehat{COB}$, cả hai góc này đều nhỏ hơn góc vuông nên là 2 góc nhọn.
\textbf{Đáp án của câu hỏi là A}
}
\end{ex}

\begin{ex} %[3H1V18]
Trong một ngày đêm (24 giờ), kim giờ và kim phút của đồng hồ tạo thành góc vuông đúng bao nhiêu lần vào các giờ đúng?
\choice
{$2$ lần}
{$12$ lần}
{\True $4$ lần}
{$24$ lần}
\loigiai{
Vào các giờ đúng, góc vuông chỉ xuất hiện vào lúc 3 giờ và 9 giờ. Một ngày đêm có hai lượt 3 giờ (3 giờ sáng, 15 giờ) và hai lượt 9 giờ (9 giờ sáng, 21 giờ), tổng cộng có 4 lần.
\textbf{Đáp án của câu hỏi là C}
}
\end{ex}

\textbf{Phần 2. Trắc nghiệm đúng sai}
\setcounterpageref{ex}{0}

\begin{ex} %[3H1V18]
Xét tính đúng sai của các khẳng định về góc:
\choiceTF
{\True Góc vuông là góc bằng góc vuông của thước ê-ke.}
{Một hình tứ giác có thể có 5 góc vuông.}
{\True Góc bé hơn góc vuông được gọi là góc nhọn.}
{Lúc 6 giờ đúng, hai kim đồng hồ tạo thành một góc vuông.}
\loigiai{
\begin{itemize}
    \item Ý a) Đúng: Định nghĩa nhận biết góc vuông bằng ê-ke.
    \item Ý b) Sai: Tứ giác chỉ có đúng 4 đỉnh, nhiều nhất 4 góc vuông.
    \item Ý c) Đúng: Kiến thức hình học cơ bản.
    \item Ý d) Sai: Lúc 6 giờ hai kim thẳng hàng tạo thành góc bẹt ($180^\circ$).
\end{itemize}
\begin{itemize}
    \item a) \textbf{Đúng}.
    \item b) \textbf{Sai}.
    \item c) \textbf{Đúng}.
    \item d) \textbf{Sai}.
\end{itemize}
}
\end{ex}

\begin{ex} %[3H1V18]
Cho hình chữ nhật $ABCD$ có điểm $M$ nằm trên cạnh $CD$:
\choiceTF
{\True Tứ giác $ABCD$ có 4 góc vuông.}
{\True Góc đỉnh $B$, cạnh $BA$ và $BC$ là góc vuông.}
{Góc đỉnh $M$, cạnh $MA$ và $MB$ luôn luôn là góc vuông.}
{\True Góc đỉnh $A$, cạnh $AB$ và $AD$ là góc vuông.}
\loigiai{
\begin{itemize}
    \item Ý a) Đúng: Đặc điểm của hình chữ nhật.
    \item Ý b) Đúng: Góc đỉnh của hình chữ nhật.
    \item Ý c) Sai: Không nhất thiết phải vuông, phụ thuộc vị trí điểm $M$.
    \item Ý d) Đúng: Góc đỉnh của hình chữ nhật.
\end{itemize}
\begin{itemize}
    \item a) \textbf{Đúng}.
    \item b) \textbf{Đúng}.
    \item c) \textbf{Sai}.
    \item d) \textbf{Đúng}.
\end{itemize}
}
\end{ex}

\textbf{Phần 3. Trả lời ngắn}
\setcounterpageref{ex}{0}

\begin{ex} %[3H1V18]
Một hình vuông có bao nhiêu góc vuông?
\shortans{4}
\loigiai{
\textbf{Đáp số:} $4$\\
Hình vuông có 4 góc vuông ở 4 đỉnh.
}
\end{ex}

\begin{ex} %[3H1V18]
Một hình tam giác vuông có bao nhiêu góc không vuông?
\shortans{2}
\loigiai{
\textbf{Đáp số:} $2$\\
Tam giác có 3 góc, trong đó có 1 góc vuông thì còn lại $3 - 1 = 2$ góc không vuông (góc nhọn).
}
\end{ex}

\begin{ex} %[3H1V18]
Từ 1 giờ đúng đến 5 giờ đúng, có bao nhiêu lần kim giờ và kim phút tạo thành một góc vuông vào đúng giờ tròn?
\shortans{1}
\loigiai{
\textbf{Đáp số:} $1$\\
Chỉ có duy nhất lúc 3 giờ đúng tạo thành góc vuông.
}
\end{ex}

\begin{ex} %[3H1V18]
Đếm số góc vuông đỉnh $O$ được tạo bởi 4 tia xuất phát từ $O$ đôi một vuông góc (tạo thành 2 đường thẳng vuông góc cắt nhau).
\shortans{4}
\loigiai{
\textbf{Đáp số:} $4$\\
Hai đường thẳng vuông góc cắt nhau tại $O$ tạo thành đúng 4 góc vuông.
}
\end{ex}

\begin{ex} %[3H1V18]
Một tờ giấy hình chữ nhật có 4 góc vuông. Dùng kéo cắt đứt một góc chéo của tờ giấy đó. Hỏi mảnh giấy lớn còn lại có bao nhiêu góc?
\shortans{5}
\loigiai{
\textbf{Đáp số:} $5$\\
Cắt vát một góc sẽ làm mất 1 góc nhưng tạo thêm 2 góc mới: $4 - 1 + 2 = 5$ góc (hình ngũ giác).
}
\end{ex}

\begin{ex} %[3H1V18]
Có bao nhiêu chữ cái in hoa trong các chữ cái sau chứa ít nhất một góc vuông: E, F, H, L, T, O, V?
\shortans{5}
\loigiai{
\textbf{Đáp số:} $5$\\
Các chữ E, F, H, L, T đều chứa các đoạn thẳng vuông góc (có góc vuông). Chữ O không có góc, chữ V là góc nhọn. Vậy có 5 chữ cái.
}
\end{ex}
